\documentclass[11pt,a4paper]{article}
\usepackage[utf8]{inputenc}
\usepackage[T1]{fontenc}
\usepackage{mathptmx}         % Times New Roman font throughout
\usepackage[top=0.75in, bottom=0.75in, left=1in, right=1in]{geometry}
\usepackage{listings}
\usepackage{needspace}
\usepackage{hyperref}

% Clean black hyperlinks
\hypersetup{
    colorlinks=true,
    linkcolor=black,
    urlcolor=black,
    citecolor=black
}

% Simple plain page style (no headers, page number at bottom center)
\pagestyle{plain}

% Simple listings setup (Times New Roman, no boxes, no line numbers)
\lstset{
    basicstyle=\fontfamily{ptm}\selectfont\normalsize,
    breaklines=true,
    breakatwhitespace=false,
    numbers=none,
    frame=none,
    showstringspaces=false,
    keepspaces=true,
    columns=fullflexible,
    tabsize=4,
    lineskip=-0.8pt
}

\begin{document}
\noindent\textbf{\large Program Name:} Design and Implement a Lexical Analyzer using C Language\par\vspace{0.6em}
\noindent\textbf{\large Program:}\par\vspace{0.3em}
\begin{lstlisting}
#include <stdio.h>
#include <stdlib.h>
#include <string.h>
#include <ctype.h>

int isKeyword(char buffer[])
{
    char keyword[32][10] = {
        "auto", "break", "case", "char", "const",
        "continue", "default", "do", "double", "else",
        "enum", "extern", "float", "for", "goto",
        "if", "int", "long", "register", "return",
        "short", "signed", "sizeof", "static", "struct",
        "switch", "typedef", "union", "unsigned", "void",
        "volatile", "while"
    };

    int i;

    for (i = 0; i < 32; i++)
    {
        if (strcmp(keyword[i], buffer) == 0)
            return 1;
    }

    return 0;
}

int main()
{
    char ch;
    char buffer[50];

    char operators[] = "+-*/%=";
    char specialch[] = ";,(){}";

    FILE *fp;
    int i, j = 0;

    fp = fopen("input.dat", "r");
    if (fp == NULL)
    {
        fp = fopen("program.txt", "r");
    }

    if (fp == NULL)
    {
        printf("Error opening file\n");
        exit(0);
    }

    while ((ch = fgetc(fp)) != EOF)
    {
        /* Check operators */
        for (i = 0; i < 6; i++)
        {
            if (ch == operators[i])
            {
                printf("%c is special character\n", ch);
                break;
            }
        }

        /* Check special characters */
        for (i = 0; i < 6; i++)
        {
            if (ch == specialch[i])
            {
                printf("%c is special character\n", ch);
                break;
            }
        }

        /* Build a word or number */
        if (isalnum(ch))
        {
            buffer[j++] = ch;
        }
        else if (j != 0)
        {
            buffer[j] = '\0';
            j = 0;

            /* Check whether token is a keyword */
            if (isKeyword(buffer))
            {
                printf("%s is keyword\n", buffer);
            }
            /* Check whether token is a number */
            else if (isdigit(buffer[0]))
            {
                int flag = 1;

                for (i = 0; buffer[i] != '\0'; i++)
                {
                    if (!isdigit(buffer[i]))
                    {
                        flag = 0;
                        break;
                    }
                }

                if (flag)
                    printf("%s is constant\n", buffer);
                else
                    printf("%s is identifier\n", buffer);
            }
            else
            {
                printf("%s is identifier\n", buffer);
            }
        }
    }

    /* Process last token if file doesn't end with whitespace */
    if (j != 0)
    {
        buffer[j] = '\0';

        if (isKeyword(buffer))
        {
            printf("%s is keyword\n", buffer);
        }
        else if (isdigit(buffer[0]))
        {
            printf("%s is constant\n", buffer);
        }
        else
        {
            printf("%s is identifier\n", buffer);
        }
    }

    fclose(fp);

    return 0;
}
\end{lstlisting}

\clearpage
\noindent\textbf{\underline{Output:}}\par\vspace{0.4em}
\noindent\textbf{input.dat:}\par\vspace{0.2em}
\begin{lstlisting}
int a=10;
float b=a+20;
if(b>30){
return b;
}
\end{lstlisting}
\vspace{0.4em}
\noindent\textbf{Output:}\par\vspace{0.2em}
\begin{lstlisting}
int is keyword
= is special character
a is identifier
; is special character
10 is constant
float is keyword
= is special character
b is identifier
+ is special character
a is identifier
; is special character
20 is constant
( is special character
if is keyword
b is identifier
) is special character
30 is constant
{ is special character
return is keyword
; is special character
b is identifier
} is special character
\end{lstlisting}

\clearpage
\noindent\textbf{\large Program Name:} Implement a Lexical Analyzer using LEX Tool\par\vspace{0.6em}
\noindent\textbf{\large Program:}\par\vspace{0.3em}
\begin{lstlisting}
%{
#include <stdio.h>
%}

%%
int|float|for|if|else|while {printf("keyword ");}
[a-zA-Z][a-zA-Z0-9]* {printf("identifier ");}
%%

int main() {
    printf("enter ip: ");
    yylex();
    return 0;
}

int yywrap() {
    return 1;
}
\end{lstlisting}

\needspace{4\baselineskip}
\vspace{0.6em}
\noindent\textbf{\underline{Output:}}\par\vspace{0.4em}
\begin{lstlisting}
enter ip: int a float b
keyword identifier keyword identifier
\end{lstlisting}

\clearpage
\noindent\textbf{\large Program Name:} Display Number of Lines, Words, Spaces, and Characters using LEX\par\vspace{0.6em}
\noindent\textbf{\large Program:}\par\vspace{0.3em}
\begin{lstlisting}
%{
#include <stdio.h>
int sc = 0, wc = 0, lc = 0, cc = 0;
%}

%%
[\n]      { lc++; cc += yyleng; }
[ \t]     { sc++; cc += yyleng; }
[^\t\n ]+ { wc++; cc += yyleng; }
%%

int main(int argc, char* argv[]) {
    printf("Enter the input:\n");
    yylex();
    printf("The number of lines=%d\n", lc);
    printf("The number of spaces=%d\n", sc);
    printf("The number of words=%d\n", wc);
    printf("The number of characters are=%d\n", cc);
    return 0;
}

int yywrap() {
    return 1;
}
\end{lstlisting}

\needspace{8\baselineskip}
\vspace{0.6em}
\noindent\textbf{\underline{Output:}}\par\vspace{0.4em}
\begin{lstlisting}
Enter the input:
Hello world
This is compiler design lab

The number of lines=2
The number of spaces=5
The number of words=7
The number of characters are=40
\end{lstlisting}

\clearpage
\noindent\textbf{\large Program Name:} Convert Substring 'abc' to 'ABC' using LEX\par\vspace{0.6em}
\noindent\textbf{\large Program:}\par\vspace{0.3em}
\begin{lstlisting}
%{
#include <stdio.h>
%}

%%
"abc" {printf("ABC");}
. {putchar(yytext[0]);}
%%

int main() {
    yylex();
    return 0;
}

int yywrap() {
    return 1;
}
\end{lstlisting}

\needspace{4\baselineskip}
\vspace{0.6em}
\noindent\textbf{\underline{Output:}}\par\vspace{0.4em}
\begin{lstlisting}
Input: welcome to abc lab
Output: welcome to ABC lab
\end{lstlisting}

\clearpage
\noindent\textbf{\large Program Name:} Count Vowels and Consonants using LEX\par\vspace{0.6em}
\noindent\textbf{\large Program:}\par\vspace{0.3em}
\begin{lstlisting}
%{
#include <stdio.h>
int v = 0;
int c = 0;
%}

%%
[aeiouAEIOU] { v++; }
[a-zA-Z]     { c++; }
[ \t\n]      ;
.            ;
%%

int main() {
    printf("Enter string: ");
    yylex();
    printf("Vowels: %d\n", v);
    printf("Consonants: %d\n", c);
    return 0;
}

int yywrap() {
    return 1;
}
\end{lstlisting}

\needspace{5\baselineskip}
\vspace{0.6em}
\noindent\textbf{\underline{Output:}}\par\vspace{0.4em}
\begin{lstlisting}
Enter string: Engineering
Vowels: 5
Consonants: 6
\end{lstlisting}

\clearpage
\noindent\textbf{\large Program Name:} Valid Arithmetic Expression using YACC\par\vspace{0.6em}
\noindent\textbf{\large Program:}\par\vspace{0.3em}
\noindent\textbf{Lexer Specification (prg6.l):}\par\vspace{0.2em}
\begin{lstlisting}
%{
#include <stdio.h>
#include "y.tab.h"
extern int yylval;
%}

%%
[0-9]+  { yylval = atoi(yytext); return NUMBER; }
[ \t]   ;
[\n]    return 0;
.       return yytext[0];
%%

int yywrap() { return 1; }
\end{lstlisting}

\noindent\textbf{YACC Parser Specification (prg6.y):}\par\vspace{0.2em}
\begin{lstlisting}
%{
#include <stdio.h>
int flag = 0;
int yylex(void);
void yyerror();
%}

%token NUMBER
%left '+' '-'
%left '*' '/' '%'
%left '(' ')'

%%
ArithmeticExpression: E { printf("\nResult=%d\n", $$); return 0; };
E: E '+' E { $$ = $1 + $3; } | E '-' E { $$ = $1 - $3; }
 | E '*' E { $$ = $1 * $3; } | E '/' E { $$ = $1 / $3; }
 | E '%' E { $$ = $1 % $3; } | '(' E ')' { $$ = $2; }
 | NUMBER  { $$ = $1; } ;
%%

void main() {
    printf("Enter an arithmetic expression: ");
    yyparse();
    if (flag == 0) printf("\nEntered Arithmetic expression is valid\n\n");
}

void yyerror() {
    printf("\nEntered Arithmetic expression is invalid\n\n");
    flag = 1;
}
\end{lstlisting}

\clearpage
\noindent\textbf{\underline{Output:}}\par\vspace{0.4em}
\begin{lstlisting}
Enter an arithmetic expression: (5+3)*2-4/2

Result=14

Entered Arithmetic expression is valid
\end{lstlisting}

\clearpage
\noindent\textbf{\large Program Name:} Valid Identifier using YACC\par\vspace{0.6em}
\noindent\textbf{\large Program:}\par\vspace{0.3em}
\noindent\textbf{Lexer Specification (prg5.l):}\par\vspace{0.2em}
\begin{lstlisting}
%{
#include <stdio.h>
#include "y.tab.h"
%}

%%
[a-zA-Z_] return letter;
[0-9]     return digit;
.         return yytext[0];
\n        return 0;
%%

int yywrap() {
    return 1;
}
\end{lstlisting}

\noindent\textbf{YACC Parser Specification (prg5.y):}\par\vspace{0.2em}
\begin{lstlisting}
%{
#include<stdio.h>
int valid=1;
int yylex(void);
int yyerror();
%}
%token digit letter
%%
start: letter s
s: letter s | digit s |;
%%
int yyerror() {
    printf("invalid identifier.\n");
    valid=0;
    return 0;
}
int main() {
    printf("enter identifier:");
    yyparse();
    if(valid)
        printf("Valid Identifier.\n");
    return 0;
}
\end{lstlisting}

\needspace{4\baselineskip}
\vspace{0.6em}
\noindent\textbf{\underline{Output:}}\par\vspace{0.4em}
\begin{lstlisting}
enter identifier: abc
Valid Identifier.
\end{lstlisting}

\clearpage
\noindent\textbf{\large Program Name:} Calculator using LEX and YACC\par\vspace{0.6em}
\noindent\textbf{\large Program:}\par\vspace{0.3em}
\noindent\textbf{Lexer Specification (calculator.l):}\par\vspace{0.2em}
\begin{lstlisting}
%{
/* Definition section */
#include<stdio.h>
#include<stdlib.h>
#include "y.tab.h"
extern int yylval;
%}

/* Rule Section */
%%
[0-9]+ {
    yylval=atoi(yytext);
    return NUMBER;
}
[ \t] ;
[\n] return 0;
. return yytext[0];
%%

int yywrap()
{
    return 1;
}
\end{lstlisting}

\noindent\textbf{YACC Parser Specification (calculator.y):}\par\vspace{0.2em}
\begin{lstlisting}
%{
/* Definition section */
#include<stdio.h>
#include<stdlib.h>
int flag=0;
int yylex(void);
void yyerror();
%}

%token NUMBER

%left '+' '-'
%left '*' '/' '%'
%left '(' ')'
/* Rule Section */
%%

ArithmeticExpression: E{
    printf("\nResult=%d\n", $$);
    return 0;
};

E:E'+'E {$$=$1+$3;}
 |E'-'E {$$=$1-$3;}
 |E'*'E {$$=$1*$3;}
 |E'/'E {$$=$1/$3;}
 |E'%'E {$$=$1%$3;}
 |'('E')' {$$=$2;}
 | NUMBER {$$=$1;}
;
%%

//driver code
void main()
{
    printf("\nEnter Any Arithmetic Expression which can have operations Addition, Subtraction, Multiplication, Division, Modulus and Round brackets:\n");

    yyparse();
    if(flag==0)
    printf("\nEntered arithmetic expression is Valid\n\n");
}

void yyerror()
{
    printf("\nEntered arithmetic expression is Invalid\n\n");
    flag=1;
}
\end{lstlisting}

\needspace{8\baselineskip}
\vspace{0.6em}
\noindent\textbf{\underline{Output:}}\par\vspace{0.4em}
\begin{lstlisting}
Enter Any Arithmetic Expression which can have operations Addition, Subtraction, Multiplication, Division, Modulus and Round brackets:
15+4*3-(10/2)

Result=22

Entered arithmetic expression is Valid
\end{lstlisting}

\clearpage
\noindent\textbf{\large Program Name:} Epsilon Closure of States in NFA\par\vspace{0.6em}
\noindent\textbf{\large Program:}\par\vspace{0.3em}
\begin{lstlisting}
#include <stdio.h>
#include <string.h>

char states[20][20];
char result[20][20];
char copy[20];

int is_present(char state[20], int n)
{
    int i;
    for (i = 0; i < n; i++)
    {
        if (strcmp(result[i], state) == 0)
            return 1;
    }
    return 0;
}

void epsilon_closure(char start[20], FILE *INPUT)
{
    char current[20];
    char state1[20], input[20], state2[20];
    int count = 0;
    int i;

    strcpy(copy, start);
    strcpy(result[count++], start);
    i = 0;
    while (i < count)
    {
        strcpy(current, result[i]);
        rewind(INPUT);
        while (fscanf(INPUT, "%s %s %s", state1, input, state2) == 3)
        {
            if (strcmp(current, state1) == 0 && strcmp(input, "e") == 0)
            {
                if (!is_present(state2, count))
                {
                    strcpy(result[count++], state2);
                }
            }
        }
        i++;
    }

    printf("\nEpsilon closure of %s = {", copy);
    for (i = 0; i < count; i++)
    {
        printf("%s", result[i]);
        if (i < count - 1)
            printf(", ");
    }
    printf("}\n");
}

int main()
{
    FILE *INPUT;
    int n, i;

    INPUT = fopen("input.dat", "r");
    if (INPUT == NULL)
    {
        printf("Unable to open input.dat\n");
        return 1;
    }

    printf("Enter the no of states: ");
    scanf("%d", &n);
    printf("Enter the states: ");
    for (i = 0; i < n; i++)
    {
        scanf("%s", states[i]);
    }
    for (i = 0; i < n; i++)
    {
        epsilon_closure(states[i], INPUT);
    }
    fclose(INPUT);
    return 0;
}
\end{lstlisting}

\needspace{8\baselineskip}
\vspace{0.6em}
\noindent\textbf{\underline{Output:}}\par\vspace{0.4em}
\noindent\textbf{input.dat:}\par\vspace{0.2em}
\begin{lstlisting}
0 e 1
1 e 2
\end{lstlisting}
\vspace{0.4em}
\noindent\textbf{Output:}\par\vspace{0.2em}
\begin{lstlisting}
Enter the no of states: 3
Enter the states: 0 1 2

Epsilon closure of 0 = {0, 1, 2}

Epsilon closure of 1 = {1, 2}

Epsilon closure of 2 = {2}
\end{lstlisting}

\clearpage
\noindent\textbf{\large Program Name:} Conversion of NFA to DFA\par\vspace{0.6em}
\noindent\textbf{\large Program:}\par\vspace{0.3em}
\begin{lstlisting}
#include <stdio.h>
#include <stdlib.h>

struct node {
    int st;
    struct node *link;
};

struct node1 {
    int nst[20];
};

void insert(int, char, int);
int findalpha(char);
void findfinalstate(void);
int insertdfastate(struct node1);
int compare(struct node1, struct node1);
void printnewstate(struct node1);

static int set[20], nostate, noalpha, s, notransition, nofinal, start, finalstate[20],r, buffer[20];
char c;
int complete = -1;
char alphabet[20];
static int eclosure[20][20] = {0};
struct node1 hash[20];
struct node *transition[20][20] = {NULL};

void main() {
    int i, j, k, m, t, n, l;
    struct node *temp;
    struct node1 newstate = {0}, tmpstate = {0};

    printf("\nEnter No of alphabets and alphabets:\n");
    scanf("%d", &noalpha);
    getchar();

    for (i = 0; i < noalpha; i++) {
        alphabet[i] = getchar();
        getchar();
    }

    printf("Enter the number of states:\n");
    scanf("%d", &nostate);

    printf("Enter the start state:\n");
    scanf("%d", &start);

    printf("Enter the number of final states:\n");
    scanf("%d", &nofinal);

    printf("Enter the final states:\n");
    for (i = 0; i < nofinal; i++) {
        scanf("%d", &finalstate[i]);
    }

    printf("Enter no of transition:\n");
    scanf("%d", &notransition);

    printf("NOTE:- [Transition is in the form-> qno alphabet qno]\n");
    printf("NOTE:- [States number must be greater than zero]\n");
    printf("\nEnter transition:\n");

    for (i = 0; i < notransition; i++) {
        scanf("%d %c%d", &r, &c, &s);
        insert(r, c, s);
    }

    for (i = 0; i < 20; i++) {
        for (j = 0; j < 20; j++) {
            hash[i].nst[j] = 0;
        }
    }

    complete = -1;
    i = -1;

    printf("\nEquivalent DFA.....\n");
    printf("Transitions of DFA\n");

    newstate.nst[start] = start;
    insertdfastate(newstate);

    while (i != complete) {
        i++;
        newstate = hash[i];

        for (k = 0; k < noalpha; k++) {
            c = 0;

            for (j = 1; j <= nostate; j++) {
                set[j] = 0;
            }

            for (j = 1; j <= nostate; j++) {
                l = newstate.nst[j];
                if (l != 0) {
                    temp = transition[l][k];
                    while (temp != NULL) {
                        if (set[temp->st] == 0) {
                            c++;
                            set[temp->st] = temp->st;
                        }
                        temp = temp->link;
                    }
                }
            }

            printf("\n");
            if (c != 0) {
                for (m = 1; m <= nostate; m++) {
                    tmpstate.nst[m] = set[m];
                }
                insertdfastate(tmpstate);
                printnewstate(newstate);
                printf("%c\t", alphabet[k]);
                printnewstate(tmpstate);
                printf("\n");
            } else {
                printnewstate(newstate);
                printf("%c\t", alphabet[k]);
                printf("NULL\n");
            }
        }
    }

    printf("\nStates of DFA:\n");
    for (i = 0; i <= complete; i++) {
        printnewstate(hash[i]);
    }

    printf("\nAlphabets:\n");
    for (i = 0; i < noalpha; i++) {
        printf("%c\t", alphabet[i]);
    }

    printf("\nStart State:\n");
    printf("q%d", start);

    printf("\nFinal states:\n");
    findfinalstate();
}

int insertdfastate(struct node1 newstate) {
    int i;
    for (i = 0; i <= complete; i++) {
        if (compare(hash[i], newstate)) {
            return 0;
        }
    }
    complete++;
    hash[complete] = newstate;
    return 1;
}

int compare(struct node1 a, struct node1 b) {
    int i;
    for (i = 1; i <= nostate; i++) {
        if (a.nst[i] != b.nst[i]) {
            return 0;
        }
    }
    return 1;
}

void insert(int r, char c, int s) {
    int j;
    struct node *temp;
    j = findalpha(c);
    if (j == 999) {
        printf("error\n");
        exit(0);
    }
    temp = (struct node *)malloc(sizeof(struct node));
    temp->st = s;
    temp->link = transition[r][j];
    transition[r][j] = temp;
}

int findalpha(char c) {
    int i;
    for (i = 0; i < noalpha; i++) {
        if (alphabet[i] == c) {
            return i;
        }
    }
    return 999;
}

void findfinalstate() {
    int i, j, k;
    for (i = 0; i <= complete; i++) {
        for (j = 1; j <= nostate; j++) {
            for (k = 0; k < nofinal; k++) {
                if (hash[i].nst[j] == finalstate[k]) {
                    printnewstate(hash[i]);
                    printf("\t");
                    j = nostate;
                    break;
                }
            }
        }
    }
}

void printnewstate(struct node1 state) {
    int j;
    printf("{");
    for (j = 1; j <= nostate; j++) {
        if (state.nst[j] != 0) {
            printf("q%d,", state.nst[j]);
        }
    }
    printf("}\t");
}
\end{lstlisting}

\clearpage
\noindent\textbf{\underline{Output:}}\par\vspace{0.4em}
\begin{lstlisting}
Enter No of alphabets and alphabets:
2
a
b
Enter the number of states:
3
Enter the start state:
1
Enter the number of final states:
1
Enter the final states:
3
Enter no of transition:
4
NOTE:- [Transition is in the form-> qno alphabet qno]
NOTE:- [States number must be greater than zero]

Enter transition:
1 a 1
1 a 2
1 b 1
2 b 3

Equivalent DFA.....
Transitions of DFA
{q1,}   a   {q1,q2,}
{q1,}   b   {q1,}
{q1,q2,}   a   {q1,q2,}
{q1,q2,}   b   {q1,q3,}
{q1,q3,}   a   {q1,q2,}
{q1,q3,}   b   {q1,}

States of DFA:
{q1,}   {q1,q2,}   {q1,q3,}

Alphabets:
a   b

Start State:
q1

Final states:
{q1,q3,}
\end{lstlisting}

\clearpage
\noindent\textbf{\large Program Name:} Minimization of DFA\par\vspace{0.6em}
\noindent\textbf{\large Program:}\par\vspace{0.3em}
\begin{lstlisting}
#include <stdio.h>
#include <stdlib.h>

#define MAX_STATES 20
#define MAX_SYMBOLS 10

int transition[MAX_STATES][MAX_SYMBOLS];
int isFinal[MAX_STATES];

int stateCount, symbolCount;
int reachable[MAX_STATES];

void dfs(int state)
{
    if (reachable[state])
        return;

    reachable[state] = 1;

    for (int i = 0; i < symbolCount; i++)
    {
        int next = transition[state][i];

        if (next != -1)
            dfs(next);
    }
}


void removeUnreachable()
{
    for (int i = 0; i < stateCount; i++)
        reachable[i] = 0;

    dfs(0);
}

int equivalent(int s1, int s2, int group[])
{

    if (isFinal[s1] != isFinal[s2])
        return 0;

    for (int i = 0; i < symbolCount; i++)
    {
        int t1 = transition[s1][i];
        int t2 = transition[s2][i];

       
        if (t1 == -1 && t2 == -1)
            continue;

  
        if (t1 == -1 || t2 == -1)
            return 0;

       
        if (group[t1] != group[t2])
            return 0;
    }

    return 1;
}

void minimizeDFA()
{
    removeUnreachable();

    int group[MAX_STATES];
    int newGroup[MAX_STATES];


    for (int i = 0; i < stateCount; i++)
    {
        if (!reachable[i])
            group[i] = -1;
        else
            group[i] = isFinal[i];
    }

    int changed;

    do
    {
        changed = 0;

        for (int i = 0; i < stateCount; i++)
            newGroup[i] = -1;

        int nextGroup = 0;

        for (int i = 0; i < stateCount; i++)
        {
            if (group[i] == -1)
                continue;

            if (newGroup[i] != -1)
                continue;

           
            newGroup[i] = nextGroup;

            for (int j = i + 1; j < stateCount; j++)
            {
                if (group[j] != -1 &&
                    newGroup[j] == -1 &&
                    equivalent(i, j, group))
                {
                    newGroup[j] = nextGroup;
                }
            }

            nextGroup++;
        }

        for (int i = 0; i < stateCount; i++)
        {
            if (group[i] != newGroup[i])
                changed = 1;

            group[i] = newGroup[i];
        }

    } while (changed);

 
    int groupCount = 0;

    for (int i = 0; i < stateCount; i++)
    {
        if (group[i] >= groupCount)
            groupCount = group[i] + 1;
    }

    printf("\nMinimized DFA:\n");
    printf("Total minimized states: %d\n", groupCount);

    for (int i = 0; i < groupCount; i++)
    {
        int final = 0;

        for (int j = 0; j < stateCount; j++)
        {
            if (group[j] == i && isFinal[j])
            {
                final = 1;
                break;
            }
        }

        printf("\nState %d", i);

        if (final)
            printf(" (Final)");

        printf(":\n");

       
        int representative = -1;

        for (int j = 0; j < stateCount; j++)
        {
            if (group[j] == i)
            {
                representative = j;
                break;
            }
        }

        if (representative != -1)
        {
            for (int s = 0; s < symbolCount; s++)
            {
                int target = transition[representative][s];

                if (target != -1)
                {
                    printf("  On symbol %d -> State %d\n",
                           s, group[target]);
                }
            }
        }
    }

    printf("\n");
}

int main()
{
    printf("Enter the number of states: ");
    scanf("%d", &stateCount);

    if (stateCount <= 0 || stateCount > MAX_STATES)
    {
        printf("Invalid number of states.\n");
        return 1;
    }

    printf("Enter number of symbols: ");
    scanf("%d", &symbolCount);

    if (symbolCount <= 0 || symbolCount > MAX_SYMBOLS)
    {
        printf("Invalid number of symbols.\n");
        return 1;
    }

    printf("Enter the transition (-1 for no transition):\n");

    for (int i = 0; i < stateCount; i++)
    {
        for (int j = 0; j < symbolCount; j++)
        {
            printf("Transition from state %d on symbol %d: ",
                   i, j);

            scanf("%d", &transition[i][j]);

            if (transition[i][j] < -1 ||
                transition[i][j] >= stateCount)
            {
                printf("Invalid transition.\n");
                return 1;
            }
        }
    }

    for (int i = 0; i < stateCount; i++)
        isFinal[i] = 0;

    int numFinalStates;

    printf("Enter the no. of final states: ");
    scanf("%d", &numFinalStates);

    if (numFinalStates < 0 || numFinalStates > stateCount)
    {
        printf("Invalid number of final states.\n");
        return 1;
    }

    printf("Enter the final states:\n");

    for (int i = 0; i < numFinalStates; i++)
    {
        int f;
        scanf("%d", &f);

        if (f < 0 || f >= stateCount)
        {
            printf("Invalid final state.\n");
            return 1;
        }

        isFinal[f] = 1;
    }

    minimizeDFA();
    return 0;
}
\end{lstlisting}

\clearpage
\noindent\textbf{\underline{Output:}}\par\vspace{0.4em}
\begin{lstlisting}
Enter the number of states: 5
Enter number of symbols: 2
Enter the transition (-1 for no transition):
Transition from state 0 on symbol 0: 1
Transition from state 0 on symbol 1: 2
Transition from state 1 on symbol 0: 1
Transition from state 1 on symbol 1: 3
Transition from state 2 on symbol 0: 1
Transition from state 2 on symbol 1: 2
Transition from state 3 on symbol 0: 1
Transition from state 3 on symbol 1: 4
Transition from state 4 on symbol 0: 1
Transition from state 4 on symbol 1: 2
Enter the no. of final states: 1
Enter the final states:
4

Minimized DFA:
Total minimized states: 4

State 0:
  On symbol 0 -> State 1
  On symbol 1 -> State 0

State 1:
  On symbol 0 -> State 1
  On symbol 1 -> State 2

State 2:
  On symbol 0 -> State 1
  On symbol 1 -> State 3

State 3 (Final):
  On symbol 0 -> State 1
  On symbol 1 -> State 0
\end{lstlisting}

\clearpage
\noindent\textbf{\large Program Name:} Computation of FIRST and FOLLOW Sets\par\vspace{0.6em}
\noindent\textbf{\large Program:}\par\vspace{0.3em}
\begin{lstlisting}
#include <stdio.h>
#include <math.h>
#include <string.h>
#include <ctype.h>
#include <stdlib.h>

int n, m = 0, p, i = 0, j = 0;
char a[10][10], f[10];
void follow(char c);
void first(char c);

int main()
{
    int i, z;
    char c, ch;
    printf("Enter the no of productions:\n");
    scanf("%d", &n);
    printf("Enter the  productions:\n");
    for (i = 0; i < n; i++)
        scanf("%s%c", a[i], &ch);
    do
    {
        m = 0;
        printf("Enter the elements whose first and follow is to be found:");
        scanf("%c", &c);
        first(c);
        printf("First(%c)={", c);
        for (i = 0; i < m; i++)
            printf("%c", f[i]);
        printf("}\n");
        strcpy(f, " ");
        m = 0;
        follow(c);
        printf("Follow(%c)={", c);
        for (i = 0; i < m; i++)
            printf("%c", f[i]);
        printf("}\n");
        printf("Continue(0/1)?");
        scanf("%d%c", &z, &ch);
    } while (z == 1);
    return (0);
}

void first(char c)
{
    int k;
    if (!isupper(c))
        f[m++] = c;
    for (k = 0; k < n; k++)
    {
        if (a[k][0] == c)
        {
            if (a[k][2] == '#')
                f[m++] = '#';
            else if (islower(a[k][2]))
                f[m++] = a[k][2];
            else
                first(a[k][2]);
        }
    }
}

void follow(char c)
{
    if (a[0][0] == c)
        f[m++] = '$';
    for (i = 0; i < n; i++)
    {
        for (j = 2; j < strlen(a[i]); j++)
        {
            if (a[i][j] == c)
            {
                if (a[i][j + 1] != '\0')
                    first(a[i][j + 1]);
                if (a[i][j + 1] == '\0' && c != a[i][0])
                    follow(a[i][0]);
            }
        }
    }
}
\end{lstlisting}

\clearpage
\noindent\textbf{\underline{Output:}}\par\vspace{0.4em}
\begin{lstlisting}
Enter the no of productions:
3
Enter the  productions:
S=Ab
A=a
A=#
Enter the elements whose first and follow is to be found: S
First(S)={a#}
Follow(S)={$}
Continue(0/1)? 1
Enter the elements whose first and follow is to be found: A
First(A)={a#}
Follow(A)={b}
Continue(0/1)? 0
\end{lstlisting}

\clearpage
\noindent\textbf{\large Program Name:} Recursive Descent Parser\par\vspace{0.6em}
\noindent\textbf{\large Program:}\par\vspace{0.3em}
\begin{lstlisting}
#include <stdio.h>
#include <string.h>
#include <ctype.h>

char input[20];
int i = 0, error = 0;
void E(), T(), Eprime(), Tprime(), F();

int main() {
    printf("Enter an arithmetic expression: ");
    gets(input);
    E();
    if (strlen(input) == i && error == 0)
        printf("\nAccepted\n");
    else
        printf("\nRejected\n");
    return 0;
}

void E() { T(); Eprime(); }

void Eprime() {
    if (input[i] == '+') { i++; T(); Eprime(); }
}

void T() { F(); Tprime(); }

void Tprime() {
    if (input[i] == '*') { i++; Tprime(); }
}

void F() {
    if (isalnum(input[i]))
        i++;
    else if (input[i] == '(') {
        i++; E();
        if (input[i] == ')') i++;
        else error = 1;
    } else error = 1;
}
\end{lstlisting}

\clearpage
\noindent\textbf{\underline{Output:}}\par\vspace{0.4em}
\begin{lstlisting}
Enter an arithmetic expression: a+b
Accepted
\end{lstlisting}

\clearpage
\noindent\textbf{\large Program Name:} Shift Reduce Parser\par\vspace{0.6em}
\noindent\textbf{\large Program:}\par\vspace{0.3em}
\begin{lstlisting}
#include <stdio.h>
#include <string.h>

struct ProductionRule
{
    char left[10];
    char right[10];
};

int main()
{
    char input[20], stack[50], temp[50], ch[2], *token1, *token2, *substring;
    int i, j, stack_length, substring_length, stack_top, rule_count = 0;
    struct ProductionRule rules[10];

    stack[0] = '\0';

    printf("\nEnter the number of production rules:");
    scanf("%d", &rule_count);

    printf("\nEnter the number of production rules(in the form of left->right):\n");
    for (i = 0; i < rule_count; i++)
    {
        scanf("%s", temp);
        token1 = strtok(temp, "->");
        token2 = strtok(NULL, "->");
        strcpy(rules[i].left, token1);
        strcpy(rules[i].right, token2);
    }

    printf("\nEnter the input string:");
    scanf("%s", input);

    i = 0;
    while (1)
    {
        if (i < strlen(input))
        {
            ch[0] = input[i];
            ch[1] = '\0';
            i++;
            strcat(stack, ch);
            printf("%s\t", stack);
            for (int k = i; k < strlen(input); k++)
            {
                printf("%c", input[k]);
            }
            printf("\tShift %s\n", ch);
        }

        for (j = 0; j < rule_count; j++)
        {
            substring = strstr(stack, rules[j].right);
            if (substring != NULL)
            {
                stack_length = strlen(stack);
                substring_length = strlen(substring);
                stack_top = stack_length - substring_length;
                stack[stack_top] = '\0';
                strcat(stack, rules[j].left);
                printf("%s\t", stack);
                for (int k = i; k < strlen(input); k++)
                {
                    printf("%c", input[k]);
                }
                printf("\tReduce %s-> %s\n", rules[j].left, rules[j].right);
                j = -1;
            }
        }

        if (strcmp(stack, rules[0].left) == 0 && i == strlen(input))
        {
            printf("\nAccepted\n");
            break;
        }

        if (i == strlen(input))
        {
            printf("\nNot Accepted\n");
            break;
        }
    }

    return 0;
}
\end{lstlisting}

\needspace{8\baselineskip}
\vspace{0.6em}
\noindent\textbf{\underline{Output:}}\par\vspace{0.4em}
\begin{lstlisting}
Enter the number of production rules: 3
Enter the number of production rules(in the form of left->right):
E->E+E
E->E*E
E->a

Enter the input string: a+a*a
\end{lstlisting}
\vspace{0.2em}
{\setlength{\tabcolsep}{16pt}
\noindent\begin{tabular}{@{}lll@{}}
a & +a*a & Shift a \\
E & +a*a & Reduce E-> a \\
E+ & a*a & Shift + \\
E+a & *a & Shift a \\
E+E & *a & Reduce E-> a \\
E & *a & Reduce E-> E+E \\
E* & a & Shift * \\
E*a & & Shift a \\
E*E & & Reduce E-> a \\
E & & Reduce E-> E*E \\
\end{tabular}}
\vspace{0.4em}
\begin{lstlisting}
Accepted
\end{lstlisting}

\clearpage
\noindent\textbf{\large Program Name:} Constant Propagation (Code Optimization)\par\vspace{0.6em}
\noindent\textbf{\large Program:}\par\vspace{0.3em}
\begin{lstlisting}
#include <stdio.h>
#include <string.h>
#include <stdlib.h>
#include <ctype.h>

void input();
void output();
void change(int p,char *res);
void constant();

struct expr
{
    char op[2],op1[5],op2[5],res[5];
    int flag;
}arr[10];

int n;

void main()
{
    input();
    constant();
    output();
}

void input()
{
    int i;
    printf("\n\nEnter the number of expressions: ");
    scanf("%d",&n);
    printf("Enter the input: \n");
    for(i=0;i<n;i++)
    {
        scanf("%s%s%s%s",arr[i].op,arr[i].op1,arr[i].op2,arr[i].res);
        arr[i].flag=0;
    }
}

void constant()
{
    int i;
    int op1,op2,res;
    char op,res1[5];
    for(i=0;i<n;i++)
    {
        if(isdigit(arr[i].op1[0]) && isdigit(arr[i].op2[0]) || strcmp(arr[i].op,"=")==0)
        {
            op1=atoi(arr[i].op1);
            op2=atoi(arr[i].op2);
            op=arr[i].op[0];
            switch(op)
            {
                case '+': res=op1+op2; break;
                case '-': res=op1-op2; break;
                case '*': res=op1*op2; break;
                case '/': res=op1/op2; break;
                case '=': res=op1; break;
            }
            sprintf(res1,"%d",res);
            arr[i].flag=1;
            change(i,res1);
        }
    }
}

void output()
{
    int i=0;
    printf("\nThe optimized code is:\n");
    for(i=0;i<n;i++)
    {
        if(!arr[i].flag)
        {
            printf("%s %s %s %s\n",arr[i].op,arr[i].op1,arr[i].op2,arr[i].res);
        }
    }
}

void change(int p,char *res)
{
    int i;
    for(i=p+1;i<n;i++)
    {
        if(strcmp(arr[p].res,arr[i].op1)==0)
        {
            strcpy(arr[i].op1,res);
        }
        else if(strcmp(arr[p].res,arr[i].op2)==0)
        {
            strcpy(arr[i].op2,res);
        }
    }
}
\end{lstlisting}

\clearpage
\noindent\textbf{\underline{Output:}}\par\vspace{0.4em}
\begin{lstlisting}
Enter the number of expressions: 3

Enter the input:

= 3 - a

+ a t1 t2

* t2 a t3

The optimized code is:

+ 3 t1 t2

* t2 3 t3
\end{lstlisting}

\clearpage
\noindent\textbf{\large Program Name:} Intermediate Code Generation (Three-Address Code)\par\vspace{0.6em}
\noindent\textbf{\large Program:}\par\vspace{0.3em}
\begin{lstlisting}
#include <stdio.h>
#include <string.h>
#include <ctype.h>

int isp(char item);
void output(char item);
void push(char item);
char pop(void);
void quad(void);

char exp[20];
char res[20];
char a[20];
int top = 0, z = 0, i = 0, l;
char x, item;

void main()
{
    printf("Enter the infix expression: ");
    gets(exp);
    l = strlen(exp);
    push('#');

    while ((item = exp[i]) != '\0')
    {
        if (isalpha(exp[i]))
            output(item);
        else if (item == '+' || item == '-' || item == '*' || item == '/' || item == '^')
            push(item);
        else if (item == '(')
            push(item);
        else if (item == ')')
        {
            while ((x = pop()) != '(')
                output(x);
        }
        else if (isp(x = pop()) < isp(item))
        {
            push(x);
            push(item);
        }
        else
        {
            output(x);
            push(item);
        }
        i++;
    }

    while ((x = pop()) != '#')
        output(x);

    printf("Postfix expression: ");
    puts(res);
    quad();
}

int isp(char item)
{
    if ((item == '+') || (item == '-'))
        return 1;
    else if ((item == '*') || (item == '/'))
        return 2;
    else if ((item == '^'))
        return 3;
    else
        return 0;
}

void output(char item)
{
    res[z++] = item;
}

void push(char item)
{
    a[++top] = item;
}

char pop(void)
{
    item = a[top--];
    return item;
}

void quad()
{
    int i, x = 0;
    char m, n, temp, str1[5], str2[5];
    printf("\nOperator\top1\top2\tresult\n");
    printf("---------------------------------");

    for (i = 0; i < z; i++)
    {
        if (isalnum(res[i]))
        {
            push(res[i]);
        }
        else
        {
            if (isalpha(m = pop()))
            {
                str1[0] = m;
                str1[1] = '\0';
            }
            else
            {
                str1[0] = 't';
                str1[1] = m;
                str1[2] = '\0';
            }

            if (isalpha(n = pop()))
            {
                str2[0] = n;
                str2[1] = '\0';
            }
            else
            {
                str2[0] = 't';
                str2[1] = n;
                str2[2] = '\0';
            }

            x++;
            printf("\n%-8c\t%-7s\t%-7s\tt%d\n", res[i], str2, str1, x);
            temp = x + '0';
            push(temp);
        }
    }
}
\end{lstlisting}

\needspace{8\baselineskip}
\vspace{0.6em}
\noindent\textbf{\underline{Output:}}\par\vspace{0.4em}
\begin{lstlisting}
Enter the infix expression: a+(b*c)-d
Postfix expression: abc*d-+
\end{lstlisting}
\vspace{0.3em}
{\setlength{\tabcolsep}{22pt}
\noindent\begin{tabular}{@{}llll@{}}
Operator & op1 & op2 & result \\[0.1em]
\multicolumn{4}{@{}l@{}}{------------------------------------------------------------} \\
{}* & b & c & t1 \\[0.8em]
-- & t1 & d & t2 \\[0.8em]
+ & a & t2 & t3 \\
\end{tabular}}

\end{document}
